\documentclass[12pt]{article}

\usepackage[a4paper,margin=1in]{geometry}
\usepackage{amsmath, amssymb}
\usepackage{enumitem}
\usepackage{array}
\usepackage{titlesec}

\titleformat{\section}{\large\bfseries}{\thesection}{1em}{}

\title{
    \textbf{Madrasa-e-Riyaziyat 2026}\\
    \large Lesson Plan
}

\author{}
\date{}

\begin{document}

\maketitle

\section{Modules}

This years madarsa has two main modules:

\begin{itemize}
    \item \textbf{Graph Theory Module}: Core mathematical foundations and structured problem solving related to graph theory
    \item \textbf{Miscellaneous Topics Module}: Various mathematics, puzzles, simulations, and other fun concepts.
\end{itemize}

\section{Daily Schedule (tentative)}

Each day is divided into two main sessions:

\begin{center}
\begin{tabular}{|c|c|}
\hline
\textbf{Time} & \textbf{Session} \\
\hline
9:00 -- 11:00 & Graph Theory \\
\hline
11:00 -- 1:00 & Break \\
\hline
1:00 -- 3:00 & Miscellaneous Topics \\
\hline
\end{tabular}
\end{center}

\section{Graph Theory Module}

This module builds foundational understanding of graph theory through problems and applications.

\subsection*{Topics}

\begin{enumerate}[label=\arabic*.]
    \item Logic and Proofs (Introduction)
    \item Bridge of Konigsberg
    \item Art Gallery Problem
    \item 3 Utilities Problem
    \item Map Coloring Problem
    \item Graph Reconstruction Problem
    \item Moore's Problem / Hall's Marriage Theorem / Handshaking Lemma
\end{enumerate}

\section{Miscellaneous Topics Module}

This module explores mathematical ideas through puzzles, simulations, and interdisciplinary topics.

\subsection*{Topics List}

\begin{itemize}
    \item Tilings, Tessellations, Escher Patterns, Symmetry, Fibonacci Sequence
    \item Chaos Theory (Double Pendulum Simulation)
    \item Galton Board (Normal Distribution Intuition)
    \item Encoding Systems: Morse Code, Barcodes, UPC
    \item Polygon Triangulation Puzzle
    \item Monty Hall Problem Simulation
    \item Sherlock Hat Puzzle
    \item Minimum Surface Ice Box Problem
    \item Sliding Puzzle Graph Representation
    \item Rubik's Cube and Combinatorics
    \item Automata Theory (Introduction)
    \item Knights Tour Problem
    \item Counting Techniques
    \item Number Theory Puzzles
    \item Godel's Incompleteness Theorem (Introductory Idea)
    \item Knot Theory and Topology (Basic Intuition)
\end{itemize}

% \section{Teaching Philosophy}

% The goal of this program is to:
% \begin{itemize}
%     \item Develop mathematical intuition rather than memorization.
%     \item Connect abstract graph theory to real-world and visual problems.
%     \item Encourage exploration through puzzles and simulations.
%     \item Build collaborative problem-solving skills.
% \end{itemize}

\section{Notes}

\begin{itemize}
    \item Topics may be adjusted based on student engagement and pacing.
    \item Graph theory module topics are confirmed; however the miscellaneous topics module has tentative topics, from which people can pick out or add their required topics. 
    \item Miscellaneous topics can include live demonstrations or coding simulations if the lecturer wants,
    \item TThere will be instructors and educators coming in from IBA, LUMS and other universities to deliver lectures. 
\end{itemize}

\end{document}